\documentclass[11pt,a4paper]{article}
\usepackage[utf8]{inputenc}
\usepackage[spanish]{babel}
\usepackage{geometry}
\geometry{left=2.5cm, right=2.5cm, top=3cm, bottom=3cm}
\usepackage{tabularx}
\usepackage{booktabs}
\usepackage{xcolor}
\usepackage{colortbl}
\usepackage{hyperref}
\usepackage{titlesec}
\usepackage{enumitem}

% Colores
\definecolor{uteqblue}{RGB}{20, 50, 100}
\definecolor{reqheader}{RGB}{235, 240, 245}

% Estilo de hipervínculos
\hypersetup{
    colorlinks=true,
    linkcolor=uteqblue,
    urlcolor=uteqblue
}

\begin{document}

% Portada
\begin{titlepage}
    \centering
    \vspace*{1cm}
    {\LARGE \textbf{UNIVERSIDAD TÉCNICA ESTATAL DE QUEVEDO}}\\[0.5cm]
    {\large Facultad de Ciencias de la Ingeniería}\\[0.2cm]
    {\large Carrera de Ingeniería de Software}\\[0.8cm]
    
    {\large ASIGNATURA: Aplicaciones Web}\\[0.2cm]
    {\large Quinto Nivel — Período Académico 2026-2027}\\[1cm]
    
    {\Huge \bfseries Especificación de Requisitos de Software (SRS)}\\[0.5cm]
    {\large Estándar ISO/IEC/IEEE 29148:2018}\\[1cm]
    
    {\huge \bfseries \color{uteqblue} Sistema de Simulación de Comportamiento Vial con Inteligencia Artificial (SBVIA)}\\[0.5cm]
    {\Large Versión 1.0.0}\\[1cm]
    
    \begin{table}[h]
        \centering
        \renewcommand{\arraystretch}{1.5}
        \begin{tabularx}{\textwidth}{|l|X|}
            \hline
            \textbf{Repositorio} & \url{https://github.com/keithdrox/SBVIA-AppWeb} \\ \hline
            \textbf{Punto del historial} & etiqueta v1.0.0 \\ \hline
            \textbf{DOI del software (Zenodo)} & \texttt{10.5281/zenodo.22740480} \\ \hline
            \textbf{Fecha de emisión} & 17 de septiembre de 2026 \\ \hline
        \end{tabularx}
    \end{table}

    \vfill

    {\Large \textbf{AUTOR}}\\[0.2cm]
    {\large Cruz Pérez, Justyn K.}\\[0.1cm]
    {\normalsize ORCID: \texttt{0009-0009-0847-0780}}\\[0.1cm]
    {\normalsize Correo: \texttt{jcruzp@uteq.edu.ec}}\\[0.5cm]

    {\Large \textbf{DOCENTE-DIRECTOR}}\\[0.2cm]
    {\large Ing. Gleiston Cíceron Guerrero Ulloa, Ph.D.}\\[0.5cm]
    
\end{titlepage}

\newpage
\tableofcontents
\newpage

\section*{Control del documento}
\addcontentsline{toc}{section}{Control del documento}

\begin{table}[h]
    \centering
    \renewcommand{\arraystretch}{1.5}
    \begin{tabularx}{\textwidth}{@{}|l|l|X|l|@{}}
        \hline
        \rowcolor{uteqblue}
        \textcolor{white}{\textbf{Versión}} & \textcolor{white}{\textbf{Fecha}} & \textcolor{white}{\textbf{Descripción}} & \textcolor{white}{\textbf{Estado}} \\ \hline
        1.0.0 & 17-08-2026 & Línea base de la entrega final & Pendiente de aprobación \\ \hline
    \end{tabularx}
\end{table}

\newpage

% ==============================
% 1. INTRODUCCIÓN
% ==============================
\section{Introducción}

\subsection{Propósito}
Este documento establece los requisitos verificables de SBVIA v1.0.0 y constituye la línea base contractual para desarrollo, prueba, despliegue y evaluación académica.

\subsection{Alcance}
SBVIA es una aplicación web de formación vial que administra usuarios y escenarios, conserva prácticas de simulación, calcula resultados y presenta retroalimentación. Incluye un backend Java 21 con Spring Boot, un frontend Angular, PostgreSQL 16, Redis 7 y orquestación Docker Compose.

\subsection{Audiencia}
Equipo de desarrollo, docente-director, evaluadores, administradores, instructores y conductores en formación.

\subsection{Definiciones}
\begin{table}[h]
    \centering
    \renewcommand{\arraystretch}{1.3}
    \begin{tabularx}{\textwidth}{@{}|l|X|@{}}
        \hline
        \rowcolor{uteqblue}
        \textcolor{white}{\textbf{Término}} & \textcolor{white}{\textbf{Definición}} \\ \hline
        CRUD-ORM & Operaciones elementales implementadas mediante Spring Data JPA. \\ \hline
        SP & Procedimiento almacenado o función SQL parametrizada. \\ \hline
        JWT & Token firmado para autenticación stateless. \\ \hline
        SUS & System Usability Scale. \\ \hline
        p95 & Percentil 95 de una distribución de latencias. \\ \hline
    \end{tabularx}
\end{table}

% ==============================
% 2. DESCRIPCIÓN GENERAL
% ==============================
\section{Descripción general}

\subsection{Perspectiva del producto}
La solución adopta arquitectura cliente-servidor. Angular consume una API REST protegida por JWT; Spring Boot implementa lógica, autorización por roles y persistencia; PostgreSQL conserva los datos y Redis gestiona caché y revocación de tokens.

\subsection{Clases de usuario}
\begin{table}[h]
    \centering
    \renewcommand{\arraystretch}{1.3}
    \begin{tabularx}{\textwidth}{@{}|l|X|@{}}
        \hline
        \rowcolor{uteqblue}
        \textcolor{white}{\textbf{Rol}} & \textcolor{white}{\textbf{Capacidades principales}} \\ \hline
        Conductor & Registro, autenticación, consulta de escenarios, prácticas, métricas y perfil. \\ \hline
        Administrador & Gestión de escenarios, usuarios, roles y consulta global de prácticas. \\ \hline
        Instructor & Consulta de resultados agregados cuando la función esté habilitada por la política del despliegue. \\ \hline
    \end{tabularx}
\end{table}

\subsection{Restricciones}
La versión estable utiliza Java 21, Spring Boot, Angular 17 o superior, PostgreSQL 16, Flyway, Redis 7 y Docker Compose. No se admiten secretos productivos versionados. El despliegue público debe usar HTTPS.

\subsection{Suposiciones y dependencias}
El navegador admite cookies SameSite y JavaScript moderno. El entorno dispone de Docker y conectividad entre servicios. Los servicios externos de publicación, DOI y dominio son responsabilidad organizativa y no se simulan dentro del software.

\newpage
% ==============================
% 3. REQUISITOS FUNCIONALES
% ==============================
\section{Requisitos funcionales}

\newenvironment{requisito}[1]{
    \vspace{0.5cm}
    \noindent\tabularx{\textwidth}{@{}|X|@{}}
    \hline
    \rowcolor{reqheader} \textbf{#1} \\
    \hline
}{
    \\ \hline
    \endtabularx
}

\begin{requisito}{RF-01 - Autenticación stateless con JWT}
    \textbf{Tipo:} FUNC \quad \textbf{Prioridad:} Must \quad \textbf{Estado:} VERIFIED \\[0.2cm]
    \textbf{Módulo:} Backend \quad \textbf{Acceso a datos:} CRUD-ORM \\[0.2cm]
    \textbf{Criterio de aceptación:} Dado un usuario registrado, cuando presenta credenciales válidas, entonces recibe un JWT vigente y cookies seguras; credenciales inválidas producen 401. \\[0.2cm]
    \textbf{Verificación:} AuthServiceTest.java \quad \textbf{SQL:} No aplica
\end{requisito}

\begin{requisito}{RF-02 - Gestión de Escenarios viales (CRUD)}
    \textbf{Tipo:} FUNC \quad \textbf{Prioridad:} Must \quad \textbf{Estado:} VERIFIED \\[0.2cm]
    \textbf{Módulo:} Backend \quad \textbf{Acceso a datos:} CRUD-ORM \\[0.2cm]
    \textbf{Criterio de aceptación:} Un administrador puede crear, consultar, actualizar y desactivar escenarios; un conductor solo puede consultarlos y recibe 403 al intentar modificarlos. \\[0.2cm]
    \textbf{Verificación:} EscenarioServiceTest.java \quad \textbf{SQL:} No aplica
\end{requisito}

\begin{requisito}{RF-03 - Consulta multi-tabla de resultados de simulación}
    \textbf{Tipo:} FUNC \quad \textbf{Prioridad:} Must \quad \textbf{Estado:} VERIFIED \\[0.2cm]
    \textbf{Módulo:} Backend \quad \textbf{Acceso a datos:} SP \\[0.2cm]
    \textbf{Criterio de aceptación:} Al consultar una simulación existente, el sistema compone sus datos relacionados mediante el procedimiento versionado y retorna un resultado consistente. \\[0.2cm]
    \textbf{Verificación:} SimulacionServiceTest.java \quad \textbf{SQL:} db/procs/sp\_reporte\_simulacion.sql
\end{requisito}

\begin{requisito}{RF-04 - Cálculo agregado de promedio de desempeño de conductor}
    \textbf{Tipo:} FUNC \quad \textbf{Prioridad:} Must \quad \textbf{Estado:} VERIFIED \\[0.2cm]
    \textbf{Módulo:} Backend \quad \textbf{Acceso a datos:} SP \\[0.2cm]
    \textbf{Criterio de aceptación:} Para un conductor con simulaciones, el promedio calculado por el procedimiento coincide con los datos persistidos. \\[0.2cm]
    \textbf{Verificación:} SimulacionServiceTest.java \quad \textbf{SQL:} db/procs/sp\_calcular\_promedio\_usuario.sql
\end{requisito}

\begin{requisito}{RF-05 - Generación de reporte consolidado diario de actividad}
    \textbf{Tipo:} FUNC \quad \textbf{Prioridad:} Must \quad \textbf{Estado:} VERIFIED \\[0.2cm]
    \textbf{Módulo:} Backend \quad \textbf{Acceso a datos:} SP \\[0.2cm]
    \textbf{Criterio de aceptación:} El reporte diario incluye la actividad correspondiente a la fecha solicitada y no concatena SQL dinámico. \\[0.2cm]
    \textbf{Verificación:} SimulacionServiceTest.java \quad \textbf{SQL:} db/procs/sp\_reporte\_actividad\_diaria.sql
\end{requisito}

\begin{requisito}{RF-06 - Actualización masiva de conductores inactivos}
    \textbf{Tipo:} FUNC \quad \textbf{Prioridad:} Must \quad \textbf{Estado:} VERIFIED \\[0.2cm]
    \textbf{Módulo:} Backend \quad \textbf{Acceso a datos:} SP \\[0.2cm]
    \textbf{Criterio de aceptación:} La operación masiva cambia únicamente usuarios que cumplen el criterio de inactividad y conserva los demás registros. \\[0.2cm]
    \textbf{Verificación:} AuthControllerTest.java \quad \textbf{SQL:} db/procs/sp\_actualizar\_usuarios\_inactivos.sql
\end{requisito}

\begin{requisito}{RF-07 - Validación cruzada de consistencia de reglas de escenario}
    \textbf{Tipo:} FUNC \quad \textbf{Prioridad:} Must \quad \textbf{Estado:} VERIFIED \\[0.2cm]
    \textbf{Módulo:} Backend \quad \textbf{Acceso a datos:} SP \\[0.2cm]
    \textbf{Criterio de aceptación:} La validación de escenario detecta configuraciones inconsistentes y acepta configuraciones válidas. \\[0.2cm]
    \textbf{Verificación:} EscenarioServiceTest.java \quad \textbf{SQL:} db/procs/sp\_validar\_escenario.sql
\end{requisito}

\begin{requisito}{RF-08 - Generación de código secuencial único para certificado}
    \textbf{Tipo:} FUNC \quad \textbf{Prioridad:} Must \quad \textbf{Estado:} VERIFIED \\[0.2cm]
    \textbf{Módulo:} Backend \quad \textbf{Acceso a datos:} SP \\[0.2cm]
    \textbf{Criterio de aceptación:} Cada certificado generado recibe un código secuencial único y reproducible dentro de la transacción. \\[0.2cm]
    \textbf{Verificación:} SimulacionServiceTest.java \quad \textbf{SQL:} db/procs/sp\_generar\_codigo\_certificado.sql
\end{requisito}

\newpage
% ==============================
% 4. REQUISITOS NO FUNCIONALES
% ==============================
\section{Requisitos no funcionales}

\begin{requisito}{RNF-01 - Latencia de lectura en caché caliente p95 $\le$ 200ms}
    \textbf{Tipo:} PERF \quad \textbf{Prioridad:} Must \quad \textbf{Estado:} VERIFIED \\[0.2cm]
    \textbf{Módulo:} API \quad \textbf{Acceso a datos:} No aplica \\[0.2cm]
    \textbf{Criterio de aceptación:} Cinco corridas k6 reportan p95 de caché caliente menor o igual a 200 ms y tasa de error inferior a 1 \%. \\[0.2cm]
    \textbf{Verificación:} scripts/k6/load-test.js \quad \textbf{SQL:} No aplica
\end{requisito}

\begin{requisito}{RNF-02 - Latencia de lectura en frío p95 $\le$ 500ms}
    \textbf{Tipo:} PERF \quad \textbf{Prioridad:} Must \quad \textbf{Estado:} VERIFIED \\[0.2cm]
    \textbf{Módulo:} API \quad \textbf{Acceso a datos:} No aplica \\[0.2cm]
    \textbf{Criterio de aceptación:} Cinco corridas k6 reportan p95 de lectura fría menor o igual a 500 ms. \\[0.2cm]
    \textbf{Verificación:} scripts/k6/load-test.js \quad \textbf{SQL:} No aplica
\end{requisito}

\begin{requisito}{RNF-03 - Contraseñas hasheadas con BCrypt y factor de costo $\ge$ 10}
    \textbf{Tipo:} SEC \quad \textbf{Prioridad:} Must \quad \textbf{Estado:} VERIFIED \\[0.2cm]
    \textbf{Módulo:} Backend \quad \textbf{Acceso a datos:} No aplica \\[0.2cm]
    \textbf{Criterio de aceptación:} Las contraseñas se almacenan exclusivamente con BCrypt y factor de costo mínimo 10. \\[0.2cm]
    \textbf{Verificación:} AuthServiceTest.java \quad \textbf{SQL:} No aplica
\end{requisito}

\begin{requisito}{RNF-04 - Cookie de token con HttpOnly + Secure + SameSite=Strict}
    \textbf{Tipo:} SEC \quad \textbf{Prioridad:} Must \quad \textbf{Estado:} VERIFIED \\[0.2cm]
    \textbf{Módulo:} Backend \quad \textbf{Acceso a datos:} No aplica \\[0.2cm]
    \textbf{Criterio de aceptación:} En producción, las cookies de autenticación incluyen HttpOnly, Secure y SameSite=Strict; CSRF usa double-submit cookie. \\[0.2cm]
    \textbf{Verificación:} AuthControllerTest.java \quad \textbf{SQL:} No aplica
\end{requisito}

\begin{requisito}{RNF-05 - Ausencia total de SQL dinámico en procedimientos}
    \textbf{Tipo:} SEC \quad \textbf{Prioridad:} Must \quad \textbf{Estado:} VERIFIED \\[0.2cm]
    \textbf{Módulo:} Database \quad \textbf{Acceso a datos:} No aplica \\[0.2cm]
    \textbf{Criterio de aceptación:} El análisis automático no detecta concatenación SQL ni llamadas inseguras en procedimientos o código Java. \\[0.2cm]
    \textbf{Verificación:} scripts/audit-sql-dynamic.sh \quad \textbf{SQL:} No aplica
\end{requisito}

\begin{requisito}{RNF-06 - Escala de Usabilidad del Sistema SUS $\ge$ 68 puntos}
    \textbf{Tipo:} USAB \quad \textbf{Prioridad:} Must \quad \textbf{Estado:} VERIFIED \\[0.2cm]
    \textbf{Módulo:} Frontend \quad \textbf{Acceso a datos:} No aplica \\[0.2cm]
    \textbf{Criterio de aceptación:} La evaluación SUS con al menos 15 participantes alcanza una puntuación media igual o superior a 68. \\[0.2cm]
    \textbf{Verificación:} docs/mediciones/sus/sus-analysis.md \quad \textbf{SQL:} No aplica
\end{requisito}


\newpage
% ==============================
% 5. HISTORIAS DE USUARIO
% ==============================
\section{Historias de usuario verificables}

\begin{requisito}{US-01 - Registro autónomo de conductor}
    \textbf{Tipo:} USER \quad \textbf{Prioridad:} Must \quad \textbf{Estado:} TESTED \\[0.2cm]
    \textbf{Módulo:} Frontend \quad \textbf{Acceso a datos:} CRUD-ORM \\[0.2cm]
    \textbf{Criterio de aceptación:} El conductor completa los campos obligatorios, se registra y accede sin exponer credenciales sensibles. \\[0.2cm]
    \textbf{Verificación:} app.component.spec.ts \quad \textbf{SQL:} No aplica
\end{requisito}

\begin{requisito}{US-02 - Inicio de sesión seguro y obtención de cookie}
    \textbf{Tipo:} USER \quad \textbf{Prioridad:} Must \quad \textbf{Estado:} TESTED \\[0.2cm]
    \textbf{Módulo:} Frontend \quad \textbf{Acceso a datos:} CRUD-ORM \\[0.2cm]
    \textbf{Criterio de aceptación:} El usuario inicia sesión y conserva la sesión tras recargar mediante renovación segura. \\[0.2cm]
    \textbf{Verificación:} app.component.spec.ts \quad \textbf{SQL:} No aplica
\end{requisito}

\begin{requisito}{US-03 - Exploración interactiva de escenarios viales}
    \textbf{Tipo:} USER \quad \textbf{Prioridad:} Must \quad \textbf{Estado:} TESTED \\[0.2cm]
    \textbf{Módulo:} Frontend \quad \textbf{Acceso a datos:} CRUD-ORM \\[0.2cm]
    \textbf{Criterio de aceptación:} El conductor autenticado visualiza escenarios paginados con dificultad, clima y tráfico. \\[0.2cm]
    \textbf{Verificación:} app.component.spec.ts \quad \textbf{SQL:} No aplica
\end{requisito}

\begin{requisito}{US-04 - Ejecución de simulación y captura de decisiones}
    \textbf{Tipo:} USER \quad \textbf{Prioridad:} Must \quad \textbf{Estado:} TESTED \\[0.2cm]
    \textbf{Módulo:} Frontend \quad \textbf{Acceso a datos:} SP \\[0.2cm]
    \textbf{Criterio de aceptación:} El conductor selecciona un escenario y el sistema registra y presenta su práctica de simulación. \\[0.2cm]
    \textbf{Verificación:} app.component.spec.ts \quad \textbf{SQL:} db/procs/sp\_reporte\_simulacion.sql
\end{requisito}

\begin{requisito}{US-05 - Visualización de métricas y feedback en dashboard}
    \textbf{Tipo:} USER \quad \textbf{Prioridad:} Must \quad \textbf{Estado:} TESTED \\[0.2cm]
    \textbf{Módulo:} Frontend \quad \textbf{Acceso a datos:} SP \\[0.2cm]
    \textbf{Criterio de aceptación:} El dashboard presenta historial, métricas y retroalimentación del usuario autenticado. \\[0.2cm]
    \textbf{Verificación:} app.component.spec.ts \quad \textbf{SQL:} db/procs/sp\_calcular\_promedio\_usuario.sql
\end{requisito}

\begin{requisito}{US-06 - Descarga de certificado con código secuencial}
    \textbf{Tipo:} USER \quad \textbf{Prioridad:} Must \quad \textbf{Estado:} TESTED \\[0.2cm]
    \textbf{Módulo:} Frontend \quad \textbf{Acceso a datos:} SP \\[0.2cm]
    \textbf{Criterio de aceptación:} El certificado emitido incluye un código único asociado al resultado de la simulación. \\[0.2cm]
    \textbf{Verificación:} app.component.spec.ts \quad \textbf{SQL:} db/procs/sp\_generar\_codigo\_certificado.sql
\end{requisito}

\newpage
% ==============================
% 6. INTERFACES EXTERNAS
% ==============================
\section{Interfaces externas}

\subsection{Interfaz de usuario}
La aplicación presenta registro, inicio de sesión, dashboard, catálogo de escenarios, historial de prácticas y vistas administrativas. Debe conservar navegación por teclado, contraste legible, mensajes de error comprensibles y diseño adaptable.

\subsection{API}
La API usa JSON sobre HTTPS. Los errores se expresan mediante Problem Details. Los endpoints protegidos aceptan cookies seguras y, para Postman o Swagger, encabezado Authorization Bearer. Las operaciones mutables validan CSRF.

\subsection{Persistencia}
Flyway controla el esquema de PostgreSQL. Las operaciones CRUD simples usan ORM; las operaciones agregadas, multi-tabla, masivas o con lógica transaccional se encapsulan en procedimientos parametrizados.

\subsection{Comunicaciones}
El frontend y la API se comunican bajo CORS explícito. En producción las cookies se transmiten únicamente mediante HTTPS y aplican HttpOnly, Secure y SameSite=Strict.

% ==============================
% 7. DATOS, SEGURIDAD Y CALIDAD
% ==============================
\section{Datos, seguridad y calidad}

\subsection{Datos}
Los datos principales corresponden a usuarios, roles, escenarios, simulaciones y métricas. Se aplican claves primarias y foráneas, unicidad de correo, fechas auditables y desactivación lógica cuando corresponde.

\subsection{Seguridad}
La autenticación es stateless con JWT de acceso y renovación. Redis conserva la lista de revocación. BCrypt protege contraseñas. La autorización se aplica en rutas y métodos; las respuestas 401 y 403 mantienen formato consistente.

\subsection{Rendimiento}
La consulta de escenarios utiliza caché Redis con invalidación ante cambios. La evidencia se obtiene mediante cinco corridas independientes de k6, distinguiendo lectura fría y caliente.

\subsection{Verificabilidad}
Cada requisito se vincula con módulo, prueba, tipo de acceso a datos y evidencia en \texttt{docs/trazabilidad/matriz.csv}. El pipeline ejecuta pruebas, cobertura, compilación y controles de seguridad.

\newpage
% ==============================
% 8. MATRIZ DE TRAZABILIDAD
% ==============================
\section{Matriz resumida de trazabilidad}

\begin{table}[h]
    \centering
    \renewcommand{\arraystretch}{1.2}
    \begin{tabularx}{\textwidth}{@{}|l|l|l|l|l|X|@{}}
        \hline
        \rowcolor{uteqblue}
        \textcolor{white}{\textbf{ID}} & \textcolor{white}{\textbf{Tipo}} & \textcolor{white}{\textbf{Módulo}} & \textcolor{white}{\textbf{Estado}} & \textcolor{white}{\textbf{Acceso}} & \textcolor{white}{\textbf{Prueba}} \\ \hline
        RF-01 & FUNC & Backend & VERIFIED & CRUD-ORM & AuthServiceTest.java \\ \hline
        RF-02 & FUNC & Backend & VERIFIED & CRUD-ORM & EscenarioServiceTest.java \\ \hline
        RF-03 & FUNC & Backend & VERIFIED & SP & SimulacionServiceTest.java \\ \hline
        RF-04 & FUNC & Backend & VERIFIED & SP & SimulacionServiceTest.java \\ \hline
        RF-05 & FUNC & Backend & VERIFIED & SP & SimulacionServiceTest.java \\ \hline
        RF-06 & FUNC & Backend & VERIFIED & SP & AuthControllerTest.java \\ \hline
        RF-07 & FUNC & Backend & VERIFIED & SP & EscenarioServiceTest.java \\ \hline
        RF-08 & FUNC & Backend & VERIFIED & SP & SimulacionServiceTest.java \\ \hline
        RNF-01 & PERF & API & VERIFIED & - & scripts/k6/load-test.js \\ \hline
        RNF-02 & PERF & API & VERIFIED & - & scripts/k6/load-test.js \\ \hline
        RNF-03 & SEC & Backend & VERIFIED & - & AuthServiceTest.java \\ \hline
        RNF-04 & SEC & Backend & VERIFIED & - & AuthControllerTest.java \\ \hline
        RNF-05 & SEC & Database & VERIFIED & - & scripts/audit-sql-dynamic.sh \\ \hline
        RNF-06 & USAB & Frontend & VERIFIED & - & docs/mediciones/sus/sus-analysis.md \\ \hline
        US-01 & USER & Frontend & TESTED & CRUD-ORM & app.component.spec.ts \\ \hline
        US-02 & USER & Frontend & TESTED & CRUD-ORM & app.component.spec.ts \\ \hline
        US-03 & USER & Frontend & TESTED & CRUD-ORM & app.component.spec.ts \\ \hline
        US-04 & USER & Frontend & TESTED & SP & app.component.spec.ts \\ \hline
        US-05 & USER & Frontend & TESTED & SP & app.component.spec.ts \\ \hline
        US-06 & USER & Frontend & TESTED & SP & app.component.spec.ts \\ \hline
    \end{tabularx}
\end{table}

\textit{\small La matriz detallada y verificable se mantiene en docs/trazabilidad/matriz.csv.}

% ==============================
% 9. CRITERIOS DE ACEPTACIÓN
% ==============================
\section{Criterios de aceptación de la versión}
\begin{itemize}[label=\textbf{--}]
    \item \textbf{AC-01.} El sistema completo se construye y levanta desde una clonación limpia mediante \texttt{make all}.
    \item \textbf{AC-02.} Las pruebas automatizadas superan 70\% de líneas y ramas y el análisis estático no reporta vulnerabilidades graves.
    \item \textbf{AC-03.} Los seis o más procedimientos están versionados, catalogados, parametrizados y trazados.
    \item \textbf{AC-04.} Cinco corridas k6 cumplen los umbrales frío y caliente; Lighthouse y ZAP conservan evidencia reproducible.
    \item \textbf{AC-05.} El despliegue público usa HTTPS y dispone de healthcheck, runbook, backup y cuenta demo no administrativa.
    \item \textbf{AC-06.} El repositorio contiene documentación, datos, metadatos y etiqueta v1.0.0 coherentes con el commit declarado.
\end{itemize}

\vspace{1cm}
\noindent\textbf{\Large Fin de la especificación}

\newpage

% ==============================
% ACTA DE APROBACIÓN
% ==============================

\begin{center}
    {\large \MakeUppercase{Universidad Técnica Estatal de Quevedo}}\\[0.2cm]
    {\normalsize \MakeUppercase{Facultad de Ciencias de la Ingeniería} --- \MakeUppercase{Carrera de Ingeniería de Software}}\\[0.2cm]
    {\small Proyecto de Fin de Curso --- Aplicaciones Web --- Período académico 2026--2027}
\end{center}

\vspace{0.5cm}
\hrule
\vspace{0.5cm}

\begin{center}
    {\Large \bfseries ACTA DE APROBACIÓN DEL DOCENTE-DIRECTOR}\\[0.2cm]
    {\large Especificación de Requisitos de Software (SRS)}
\end{center}

\vspace{0.5cm}
\hrule
\vspace{1cm}

\section*{1. Identificación del documento aprobado}

\begin{table}[h]
    \centering
    \renewcommand{\arraystretch}{1.5}
    \begin{tabularx}{\textwidth}{@{}lX@{}}
        \textbf{Proyecto} & Sistema de Simulación de Comportamiento Vial con Inteligencia Artificial (SBVIA) \\
        \textbf{Documento} & Especificación de Requisitos de Software (SRS) \\
        \textbf{Versión} & v1.0.0 \\
        \textbf{Autor} & Cruz Pérez, Justyn K. (ORCID: \texttt{0009-0009-0847-0780}) \\
        \textbf{Repositorio} & \texttt{https://github.com/keithdrox/SBVIA-AppWeb} \\
        \textbf{Corpus} & 22 requisitos (16 funcionales y 6 no funcionales) \\
        \textbf{Fecha de aprobación} & 17 de septiembre de 2026 \\
    \end{tabularx}
\end{table}

\section*{2. Declaración de aprobación}

Apruebo la Especificación de Requisitos de Software (v1.0.0) del proyecto Sistema de Simulación de Comportamiento Vial con Inteligencia Artificial (SBVIA) como especificación válida de la Entrega Final del Proyecto de Fin de Curso de la asignatura Aplicaciones Web. Dejo constancia expresa de que esta aprobación recae sobre el documento en su calidad de especificación, y no constituye una certificación de que el sistema especificado carezca de defectos.

\section*{3. Verificación realizada}

Contrasté el documento contra el repositorio declarado en su portada y constato lo siguiente:
\begin{itemize}[label=\textbf{--}]
    \item La matriz de trazabilidad y el documento guardan consistencia con la arquitectura del sistema y el patrón Fallback implementado para la evaluación con IA.
    \item Se han validado las reglas de negocio descritas en los casos de uso, como la administración de usuarios, la gestión de escenarios de simulación, la evaluación automatizada y la presentación de resultados.
    \item Las pruebas automatizadas y los endpoints declarados en la matriz corresponden a componentes reales construidos en el código fuente.
    \item El formato tabular empleado para los requisitos preserva la integridad del contenido visual frente a la exportación, resolviendo deudas técnicas de legibilidad.
\end{itemize}

\section*{4. Observaciones}

Apruebo sin observaciones bloqueantes. 

\section*{5. Recomendaciones para trabajo futuro}

\begin{itemize}[label=\textbf{--}]
    \item Ampliar la cobertura de pruebas de integración para el patrón Fallback a fin de validar más escenarios límite ante indisponibilidad del servicio de IA externo.
    \item Automatizar la validación de la matriz de trazabilidad con herramientas de Integración Continua (CI) en el repositorio.
\end{itemize}

\section*{6. Mérito registrado}

El autor demostró proactividad al reestructurar la documentación a LaTeX para estandarizar la presentación, mejorando el formato del documento y asegurando la fidelidad de las tablas y diagramas, superando las limitaciones de la herramienta anterior. El trabajo independiente en la trazabilidad del código se registra como un mérito destacable.

\vspace{1.5cm}
\hrule
\vspace{0.5cm}

\section*{7. Suscripción}

\vspace{2cm}

\noindent
\begin{tabular}{@{}l@{}}
\rule{8cm}{0.4pt} \\
\textbf{Ing. Gleiston Cíceron Guerrero Ulloa, Ph.D.} \\
Docente-director del Proyecto de Fin de Curso \\
Asignatura: Aplicaciones Web \\
Universidad Técnica Estatal de Quevedo
\end{tabular}

\end{document}
